\section{Localization and navigation}
\section{Localization}
\subsubsection{Odometry}
If you just want to move by 3 meters to test a whole body controller,
and you do not really need to know where on the map this is happening,
you want odometry. If you are interacting with one or few objects in the environment,
and do not need to search for them (ex. you're testing a door-opening controller),
then put QR codes on the object, and localize from vision ---
you do not need to know where these QR code are in the room (or that there even is a room apart
from the space between the robot and the object).

Falls apart in about 10 meters or 1-2 minutes or runtime.

\subsubsection{Maps}
Use someone else SLAM unless this is your explicit focus.

\subsubsection{Localization on maps}
Write or use someone else Kalman filter to bridge between slow localization updates
and fast odometry.

\subsection{Navigation}
You actually just use existing controllers,
the only difference is that the robot model is different.
However, some work better than others,
especially for underactuated platforms.
Assuming this is the case, consider the following controllers:
\begin{itemize}
		\item MPC for trajectory tracking
		\item QP with velocity bounds: TODO: test this better
\end{itemize}
